\chapter{Suspension flows: exponents, ergodicity, entropy}
\label{ch:suspension}

The abstract continuous-flow multiplicative ergodic theorem of
Chapter~\ref{chap:continuous} governs any measure-preserving $\R$-action carrying a
continuous-time cocycle. This chapter studies the single most important \emph{source} of
such flows --- the \emph{suspension} (special-flow, or mapping-torus) construction, which
manufactures a flow from a discrete base automorphism $T$ and a positive \emph{roof}
function $\tau$ --- and computes its three characteristic invariants over concrete
hyperbolic and Bernoulli bases: the Lyapunov exponents of the flow cocycle, the ergodicity
of individual time-$t$ maps, and the Kolmogorov--Sinai entropy of the flow.

The unifying theme is Abramov's principle of \emph{time rescaling}: a suspension flow runs
the base clock at speed set by the roof, so a base invariant of rate type is divided by the
mean roof $\int\tau$. This yields the Lyapunov analogue
$\lambda_{\mathrm{flow}} = \lambda_{\mathrm{base}}/\int\tau$ of Abramov's entropy formula
$h(\text{flow}) = h(\text{base})/\int\tau$ (L. M. Abramov, \emph{On the entropy of a flow},
Dokl.\ Akad.\ Nauk SSSR \textbf{128} (1959) 873--875; L. Barreira, \emph{Lyapunov Exponents},
Birkh\"auser 2017, Ch.~3, on exponents under time rescaling). The time-$1$ ergodicity holds
exactly for \emph{irrational} roofs; the time-$1$ entropy identity $h(\zeta^{(r)}_1) = H_\nu/r$
is now proved for \emph{every} roof $r > 0$ via the abstract Abramov flow-entropy homogeneity
$h(\varphi_t) = t\,h(\varphi_1)$ (\Cref{thm:flow-abramov-hom}), following Ito's elementary
generator-free argument --- retiring the former rational-roof restriction.
Throughout, the standing references for the special-flow bookkeeping are Cornfeld--Fomin--Sinai
(\emph{Ergodic Theory}, Grundlehren \textbf{245}, Springer 1982, Ch.~10--11) and the
Ambrose--Kakutani structure theorem for measurable flows (Duke Math.\ J.\ \textbf{9} (1942)
25--42).

\section{The suspension space and its invariant measure}

The suspension of a measurable automorphism $T : X \simeq X$ under a measurable roof
$\tau : X \to \R$ is the quotient of the cover $X \times \R$ by the $\Z$-action generated by
the shear $G(x, s) = (Tx,\, s - \tau x)$; a point flows upward, $\zeta_t[x, s] = [x, s + t]$,
and re-enters the base after time $\tau x$. We recall only the objects the rest of the chapter
consumes; the abstract flow theorem they feed is Chapter~\ref{chap:continuous}'s
\Cref{oseledets_flow}.

\begin{definition}[The suspension (mapping-torus) space]\label{thm:susp-space}
  \lean{ErgodicTheory.SuspensionSpace}\leanok
  The \emph{suspension space} $\Sigma = \operatorname{SuspensionSpace} T\,\tau$ is the orbit
  quotient of $X \times \R$ by the suspension $\Z$-action $G(x, s) = (Tx,\, s - \tau x)$,
  i.e.\ the mapping torus of $T$ under the roof $\tau$. As a \texttt{Quotient} it carries the
  canonical pushed-forward \texttt{MeasurableSpace}, and the projection
  $\pi = \operatorname{suspensionMk} : X \times \R \to \Sigma$ is measurable.
\end{definition}

\begin{definition}[The invariant suspension probability measure]\label{thm:susp-measure}
  \lean{ErgodicTheory.suspensionMeasure}\leanok
  \uses{thm:susp-space}
  The \emph{suspension measure} $\hat\mu = \operatorname{suspensionMeasure} T\,\tau\,\mu$ is
  obtained by restricting $\mu \times \operatorname{vol}$ to the fundamental box
  $\{(x, s) : 0 \le s < \tau x\}$, pushing forward along $\pi$, and normalising by
  $(\int\tau)^{-1}$. The generator $G$ preserves $\mu \times \operatorname{vol}$ whenever $T$
  preserves $\mu$ (\lean{ErgodicTheory.measurePreserving_suspensionGen}), because the shear is
  a fibered translation and $\R$-translation is Haar-invariant.
\end{definition}

\begin{theorem}[$\hat\mu$ is a probability measure]\label{thm:susp-prob}
  \lean{ErgodicTheory.isProbabilityMeasure_suspensionMeasure}\leanok
  \uses{thm:susp-measure}
  For a nonnegative integrable roof with $0 < \int\tau$, the normalised suspension measure
  $\hat\mu$ is a probability measure.
\end{theorem}
\begin{proof}\leanok
  \uses{thm:susp-measure}
  The raw box push-forward has total mass $\int\tau$ (Fubini: the $x$-fibre of the box is
  $[0, \tau x)$, of length $\tau x$), which is positive and finite; multiplying by
  $(\int\tau)^{-1}$ makes the total mass $1$.
\end{proof}

The upward translation descends to a genuine measure-preserving flow on $\Sigma$. Its time-$t$
map is $\operatorname{suspensionFlowMap} T\,\tau\,t$
(\lean{ErgodicTheory.suspensionFlowMap}), packaged as a \texttt{MeasurePreservingFlow}
$\operatorname{suspensionFlow}$ (\lean{ErgodicTheory.suspensionFlow},
\lean{ErgodicTheory.measurePreserving_suspensionFlowMap}), with the base cross-section
$x \mapsto [x, 0]$ recorded as $\operatorname{suspensionSection}$
(\lean{ErgodicTheory.suspensionSection}). These are exactly the ingredients
Definition~\ref{MeasurePreservingFlow} of Chapter~\ref{chap:continuous} asks for, so the
abstract flow MET applies verbatim once a flow cocycle is supplied. The rest of the chapter
supplies and analyses that cocycle.

\section{The flow Lyapunov exponent over the suspension}

A base matrix generator $A : X \to \operatorname{Mat}_d(\R)$ generates, over the cross-section
flow, a \emph{cover cocycle} on $X \times \R$: reading a base point $x$ at fibre height $h$,
advancing the special flow for time $t$ accumulates the base cocycle over the total elapsed
section time $h + t$.

\begin{definition}[The cover flow cocycle]\label{thm:susp-coverCocycle}
  \lean{ErgodicTheory.coverCocycle}\leanok
  \uses{cocycle}
  For a cover point $p = (x, h)$ the \emph{cover cocycle}
  $\operatorname{coverCocycle} A\,T\,\tau\,(x, h)\,t$ is the cross-section flow cocycle read at
  the total elapsed time $h + t$ from the base point $x$. This is the cover-level form of the
  Cornfeld--Fomin--Sinai special-flow cocycle, one step before the orbit-quotient descent.
\end{definition}

\begin{lemma}[Agreement on the base section]\label{thm:susp-coverBase}
  \lean{ErgodicTheory.coverCocycle_base}\leanok
  \uses{thm:susp-coverCocycle}
  At height $0$ the cover cocycle is the cross-section flow cocycle:
  $\operatorname{coverCocycle}(x, 0)\,t = \operatorname{flowCocycleSection} t\,x$, since the
  total elapsed time $0 + t$ collapses to $t$.
\end{lemma}
\begin{proof}\leanok
  \uses{thm:susp-coverCocycle}
  Immediate from the definition on rewriting $0 + t = t$ in the fibre coordinate.
\end{proof}

\begin{theorem}[Section multiplicativity at a return boundary]\label{thm:susp-coverReturn}
  \lean{ErgodicTheory.coverCocycle_section_returnTime}\leanok
  \uses{thm:susp-coverBase, cocycle}
  Starting on the base section at $x$, the cover cocycle over flow time
  $\operatorname{returnTime} n\,x + r$ (for $0 \le r$) factors as the residual flow over time
  $r$ from the shifted base point $T^n x$, composed on the left of the discrete base cocycle
  for the $n$ completed laps:
  \[
    \operatorname{coverCocycle}(x, 0)\bigl(\operatorname{returnTime} n\,x + r\bigr)
      = \operatorname{coverCocycle}(T^n x, 0)\,r \cdot \operatorname{cocycle} A\,T\,n\,x.
  \]
\end{theorem}
\begin{proof}\leanok
  \uses{thm:susp-coverBase}
  The accumulated matrix splits at the return boundary because the lap counter does; the return
  cocycle is multiplicative across the split, and its first $n$ laps give
  $\operatorname{cocycle} A\,T\,n\,x$ by the return identity. Reordering the split so the $n$
  completed laps sit on the right yields the stated factorisation.
\end{proof}

The \emph{matrix} cover cocycle does not descend to $\Sigma$: re-basing along $n$ orbit steps
post-multiplies it by the fixed factor $\operatorname{cocycle} A\,T\,n\,x$. What descends is the
\emph{growth rate}, because that fixed multiplicative factor becomes a fixed additive $\log$
shift, washed out by the $1/t$ Birkhoff average.

\begin{theorem}[Additive $\log$-discrepancy across an orbit step]\label{thm:susp-logDisc}
  \lean{ErgodicTheory.coverCocycle_suspensionAct_log_discrepancy}\leanok
  \uses{thm:susp-coverReturn}
  Under invertibility of the base cocycle and strict positivity of both cover-cocycle norms, for
  $\operatorname{returnTime} n\,x \le s + t$ the two $\log$-norms differ by at most a constant in
  $t$:
  \[
    \bigl|\log\norm{\operatorname{coverCocycle}(\operatorname{suspensionAct} n\,(x, s))\,t}
      - \log\norm{\operatorname{coverCocycle}(x, s)\,t}\bigr|
      \le \bigl|\log\norm{\operatorname{cocycle} A\,T\,n\,x}\bigr|
        + \bigl|\log\norm{(\operatorname{cocycle} A\,T\,n\,x)^{-1}}\bigr|.
  \]
  The absolute values keep the bound honest: an operator norm need not be $\ge 1$, so an
  individual $\log$-norm can be negative.
\end{theorem}
\begin{proof}\leanok
  \uses{thm:susp-coverReturn}
  Take $\log$ of the two operator-norm brackets supplied by \Cref{thm:susp-coverReturn} (which
  bound $\norm{\operatorname{coverCocycle}(x, s)\,t}$ between $\norm{\cdot}$ at the re-based
  point times $\norm{\operatorname{cocycle}}^{\pm1}$), using $\log$-monotonicity and
  $\log(ab) = \log a + \log b$ on the strictly positive factors, then bound each $\log$ factor by
  its absolute value.
\end{proof}

\begin{theorem}[The Lyapunov-exponent limit transfer]\label{thm:susp-limitTransfer}
  \lean{ErgodicTheory.coverCocycle_suspensionAct_tendsto_exponent}\leanok
  \uses{thm:susp-logDisc}
  If the cover-cocycle growth rate $t^{-1}\log\norm{\operatorname{coverCocycle}(x, s)\,t}$
  converges to $L$ as $t \to \infty$, the base cocycle is invertible, and both cover-cocycle
  norms are eventually strictly positive, then the growth rate at the re-based orbit point
  $\operatorname{suspensionAct}(n)\,(x, s)$ converges to the \emph{same} $L$.
\end{theorem}
\begin{proof}\leanok
  \uses{thm:susp-logDisc}
  The per-$t$ average at the re-based point lies within $C/t$ of the average at $(x, s)$, where
  $C$ is the $t$-independent additive discrepancy of \Cref{thm:susp-logDisc}; since
  $C/t \to 0$, a \texttt{Filter.Tendsto} squeeze transfers the limit.
\end{proof}

\begin{definition}[The flow Lyapunov exponent of an orbit class]\label{thm:susp-hasFlow}
  \lean{ErgodicTheory.HasFlowExponent}\leanok
  \uses{thm:susp-coverCocycle, thm:susp-space}
  $\operatorname{HasFlowExponent} q\,L$ holds when \emph{some} representative $(x, s)$ of the
  orbit class $q \in \Sigma$ carries the cover-cocycle growth rate $L$:
  $\pi(x, s) = q$ and $t^{-1}\log\norm{\operatorname{coverCocycle}(x, s)\,t} \to L$. The
  existential form lifts cleanly to the quotient; the well-definedness results below show that,
  under base-cocycle invertibility, the rate is in fact the same for every representative.
\end{definition}

\begin{theorem}[Two-sided well-definedness across a forward step]\label{thm:susp-iffForward}
  \lean{ErgodicTheory.tendsto_exponent_iff_of_suspensionAct}\leanok
  \uses{thm:susp-limitTransfer}
  If $(x_2, s_2) = \operatorname{suspensionAct}(n)\,(x, s)$ for a natural $n$, the base cocycle
  is invertible, and both cover-cocycle norms are eventually strictly positive, then the
  cover-cocycle growth rates at $(x, s)$ and at $(x_2, s_2)$ converge to one and the same $L$:
  the two \texttt{Tendsto} statements are equivalent.
\end{theorem}
\begin{proof}\leanok
  \uses{thm:susp-limitTransfer}
  The $\to$ direction is \Cref{thm:susp-limitTransfer}. The $\leftarrow$ direction is its mirror:
  the additive $\log$-discrepancy bound of \Cref{thm:susp-logDisc} is symmetric in the two points,
  so the same $1/t$ squeeze transfers the limit the other way.
\end{proof}

\begin{theorem}[Orbit-class invariance of the flow exponent]\label{thm:susp-classInv}
  \lean{ErgodicTheory.hasFlowExponent_of_suspensionAct}\leanok
  \uses{thm:susp-hasFlow}
  If two cover points are connected by a forward orbit step
  $\operatorname{suspensionAct}(n)\,(x, s) = (x_2, s_2)$, then a growth rate $L$ computed at
  $(x, s)$ witnesses $\operatorname{HasFlowExponent}(\pi(x_2, s_2))\,L$: the exponent is a
  property of the whole orbit class, read off from any representative.
\end{theorem}
\begin{proof}\leanok
  \uses{thm:susp-hasFlow}
  Orbit-equivalent points have equal $\pi$-images, so $(x, s)$ itself serves as the witness for
  the common class; its growth rate is $L$ by hypothesis. (Invertibility is needed only to
  guarantee the \emph{same} $L$ from the other representative, via \Cref{thm:susp-iffForward}.)
\end{proof}

\subsection{The Abramov exponent formula}

Disintegrating $\hat\mu$ over the base measure and feeding in the base Birkhoff limits gives the
headline: the flow exponent is the base exponent divided by the mean roof.

\begin{theorem}[The special-flow Lyapunov exponent, $\lambda_{\mathrm{flow}} = \lambda_{\mathrm{base}}/\int\tau$]\label{thm:susp-abramovExp}
  \lean{ErgodicTheory.ae_suspensionMeasure_hasFlowExponent_of_measurable}\leanok
  \uses{thm:susp-iffForward, thm:susp-classInv, thm:susp-prob}
  Under a bounded roof $c \le \tau \le C$ ($0 < c$), positive integral $0 < \int\tau$, measurable
  base generator $A$, and the base-a.e.\ Birkhoff limits --- discrete base growth rate
  $\to \lambda_{\mathrm{base}}$ and roof average $\to \int\tau$ --- for $\hat\mu$-almost every
  orbit class $q$,
  \[
    \operatorname{HasFlowExponent} q\ \Bigl(\lambda_{\mathrm{base}}\big/\!\int\tau\Bigr).
  \]
\end{theorem}
\begin{proof}\leanok
  \uses{thm:susp-iffForward, thm:susp-classInv}
  The base exponent-set measurability is supplied internally: the full-time cover-cocycle exponent
  set is rewritten pointwise as the discrete return-time exponent set (the between-returns squeeze)
  and is measurable by \texttt{measurableSet\_tendsto}, so only $A$ measurable is assumed. The
  base-a.e.\ growth and roof limits are transported to $\hat\mu$-a.e.\ orbit classes through the
  fundamental-domain disintegration, and the between-returns squeeze converts the discrete
  return-time exponent $\lambda_{\mathrm{base}}/\int\tau$ into the full-time flow exponent.
\end{proof}

\begin{theorem}[The exponent along genuine flow orbits]\label{thm:susp-abramovExpFlow}
  \lean{ErgodicTheory.ae_suspensionMeasure_hasFlowExponent_flowOrbit_of_measurable}\leanok
  \uses{thm:susp-abramovExp}
  Under the same hypotheses (and $T$ measure-preserving), for $\hat\mu$-a.e.\ $q$ there are a base
  point $x$ and a flow time $s$ with $q = \operatorname{suspensionFlow} s\,(\text{section } x)$ and
  $\operatorname{HasFlowExponent} q\,(\lambda_{\mathrm{base}}/\int\tau)$: the exponent is realised
  on the genuine measure-preserving flow orbit of a cross-section point.
\end{theorem}
\begin{proof}\leanok
  \uses{thm:susp-abramovExp}
  The disintegration already identifies each a.e.\ class with a flow-orbit point of the base
  section; carry that identification alongside the exponent conclusion of
  \Cref{thm:susp-abramovExp}.
\end{proof}

\subsection{The representative-free descent}

The predicate $\operatorname{HasFlowExponent}$ is existential over representatives. Upgrading it to
a genuine function $\Sigma \to \R$ requires closing the \emph{signed} orbit step (a backward step
$m < 0$) and paying one honest cost: a global base-cocycle invertibility hypothesis that makes the
\texttt{Quotient.lift} total.

\begin{lemma}[Strict positivity under global invertibility]\label{thm:susp-normPos}
  \lean{ErgodicTheory.coverCocycle_norm_pos}\leanok
  \uses{thm:susp-coverCocycle}
  If the base generator $A$ is everywhere invertible ($\det \ne 0$), then
  $0 < \norm{\operatorname{coverCocycle} p\,t}$ for every cover point $p$ and time $t$: the cover
  cocycle reduces to a base-cocycle iterate whose norm is strictly positive.
\end{lemma}
\begin{proof}\leanok
  \uses{thm:susp-coverCocycle}
  Unfold the cover cocycle to $\operatorname{cocycle} A\,T\,(\text{lap count})\,x$, a product of
  invertible matrices, hence nonzero with positive operator norm.
\end{proof}

\begin{theorem}[Signed-step cross-representative uniqueness]\label{thm:susp-iffSigned}
  \lean{ErgodicTheory.tendsto_exponent_iff_of_orbitRel}\leanok
  \uses{thm:susp-iffForward, thm:susp-normPos}
  If two cover points are connected by a \emph{signed} orbit step
  $\operatorname{suspensionAct} m\,(x_2, s_2) = (x_1, s_1)$ for $m \in \Z$, and the base cocycle is
  everywhere invertible, then the cover-cocycle growth rates at $(x_1, s_1)$ and $(x_2, s_2)$
  converge to the same $L$.
\end{theorem}
\begin{proof}\leanok
  \uses{thm:susp-iffForward, thm:susp-normPos}
  Both signs reduce to the forward iff of \Cref{thm:susp-iffForward} at a different base point:
  for $0 \le m$ at base $(x_2, s_2)$; for $m \le 0$ after inverting the connection to
  $\operatorname{suspensionAct}(-m)\,(x_1, s_1) = (x_2, s_2)$, at base $(x_1, s_1)$. Global
  invertibility discharges the unit-determinant and eventual strict-positivity side conditions on
  both representatives via \Cref{thm:susp-normPos}.
\end{proof}

\begin{definition}[The representative-level and descended exponents]\label{thm:susp-flowExpAt}
  \lean{ErgodicTheory.flowExponentAt}\leanok
  \uses{thm:susp-iffSigned}
  $\operatorname{repExponent} p$ (\lean{ErgodicTheory.repExponent}) is the growth-rate limit
  $\lim_t t^{-1}\log\norm{\operatorname{coverCocycle} p\,t}$ read off from a representative, with a
  fixed junk value $0$ off the convergence locus (a plain $\texttt{dite}$/$\texttt{choose}$ rather
  than $\texttt{limUnder}$, whose junk is not constant across representatives).
  $\operatorname{flowExponentAt} : \Sigma \to \R$ is its $\texttt{Quotient.lift}$, made total by
  global base-cocycle invertibility.
\end{definition}

\begin{theorem}[$\operatorname{flowExponentAt}$ reads off the exponent]\label{thm:susp-flowExpEq}
  \lean{ErgodicTheory.flowExponentAt_eq_of_hasFlowExponent}\leanok
  \uses{thm:susp-flowExpAt, thm:susp-hasFlow}
  If $q$ carries the flow exponent $L$ (some representative has cover-cocycle growth rate $L$),
  then $\operatorname{flowExponentAt} q = L$.
\end{theorem}
\begin{proof}\leanok
  \uses{thm:susp-flowExpAt}
  Well-definedness of the lift (\Cref{thm:susp-iffSigned}) transfers existence of the limit across
  any orbit step; where it exists on both sides, uniqueness of limits forces the two values to
  agree, so the descended value equals the witnessing $L$.
\end{proof}

\begin{theorem}[The representative-free Abramov exponent]\label{thm:susp-flowExpAe}
  \lean{ErgodicTheory.ae_flowExponentAt_eq_base_div_roof}\leanok
  \uses{thm:susp-flowExpEq, thm:susp-abramovExp}
  Under a bounded roof, positive $\int\tau$, measurable base generator, global invertibility
  $\forall x,\ \det(A x) \ne 0$, and the base-a.e.\ Birkhoff limits, for $\hat\mu$-almost every
  orbit class $q$ the genuine (representative-free) flow exponent equals
  $\lambda_{\mathrm{base}}/\int\tau$:
  $\operatorname{flowExponentAt} q = \lambda_{\mathrm{base}}/\int\tau$.
\end{theorem}
\begin{proof}\leanok
  \uses{thm:susp-flowExpEq, thm:susp-abramovExp}
  Combine the existential a.e.\ exponent of \Cref{thm:susp-abramovExp} with the read-off equality
  \Cref{thm:susp-flowExpEq}; the added global-invertibility hypothesis is the documented honest
  cost of upgrading the predicate to the actual lifted value.
\end{proof}

\subsection{The cat-map suspension flow}

The exponent theory is instantiated on the suspension of the Arnold cat map $\operatorname{catTorus}$
(\Cref{catTorus}) under the unit roof $\tau \equiv 1$, fed the base's \emph{own} derivative cocycle
(\Cref{derivativeCocycle}), which is the constant hyperbolic matrix
$M = \left(\begin{smallmatrix} 2 & 1 \\ 1 & 1 \end{smallmatrix}\right)$.
Its spectrum (\Cref{catTorus_constCocycle_exponents}) gives the base top exponent
$\log\bigl((3+\sqrt5)/2\bigr) > 0$, and $\int\tau = 1$, so $\lambda_{\mathrm{flow}} = \lambda_{\mathrm{base}}$.

\begin{theorem}[The cat suspension realises the base derivative-cocycle exponent]\label{thm:susp-catHasFlow}
  \lean{ErgodicTheory.CatMapToral.catSuspension_ae_hasFlowExponent}\leanok
  \uses{thm:susp-abramovExp, catTorus_constCocycle_exponents, derivativeCocycle}
  For $\hat\mu$-a.e.\ orbit class $q$ of the cat-map suspension under the unit roof, the flow
  Lyapunov exponent of the base's own derivative cocycle is
  $\operatorname{HasFlowExponent} q\ \bigl(\log((3+\sqrt5)/2)\bigr) = \lambda_{\mathrm{base}}/\int\tau$.
\end{theorem}
\begin{proof}\leanok
  \uses{thm:susp-abramovExp}
  The generator is constant in the base point, so its discrete growth rate is the deterministic
  Gelfand limit $n^{-1}\log\norm{M^n} \to \log((3+\sqrt5)/2)$ (spectral radius identified through
  the Grade-1 cat spectrum \Cref{catTorus_constCocycle_exponents}), and the unit-roof Birkhoff
  average is $1 = \int\tau$. Instantiate \Cref{thm:susp-abramovExp} at this data.
\end{proof}

\begin{theorem}[Positivity of the cat suspension exponent (issue \#30)]\label{thm:susp-catPos}
  \lean{ErgodicTheory.CatMapToral.catSuspensionFlow_ownExponent_pos}\leanok
  \uses{thm:susp-catHasFlow}
  For $\hat\mu$-a.e.\ orbit class $q$ there is a \emph{positive} $L$ with
  $\operatorname{HasFlowExponent} q\,L$: the hyperbolicity of the cat map, $(3+\sqrt5)/2 > 1$,
  makes $L = \log((3+\sqrt5)/2) > 0$.
\end{theorem}
\begin{proof}\leanok
  \uses{thm:susp-catHasFlow}
  Take $L = \log((3+\sqrt5)/2)$ from \Cref{thm:susp-catHasFlow}; positivity is
  $\log$ of a number exceeding $1$.
\end{proof}

\begin{theorem}[Abramov quotient reading, existential form]\label{thm:susp-catBaseDiv}
  \lean{ErgodicTheory.CatMapToral.catSuspension_flowExponent_eq_base_div_roof}\leanok
  \uses{thm:susp-catHasFlow, catTorus_constCocycle_exponents}
  The flow exponent equals the base top Lyapunov exponent of $M$ over the genuine ergodic cat map
  divided by the mean roof $\int\tau$; since $\int\tau = 1$ this is the same numerical value,
  exposed in its Abramov-style quotient form.
\end{theorem}
\begin{proof}\leanok
  \uses{thm:susp-catHasFlow}
  Rewrite the base top exponent as the Grade-1 spectral value $\log((3+\sqrt5)/2)$ and divide by
  $\int\tau = 1$; the result is the exponent of \Cref{thm:susp-catHasFlow}.
\end{proof}

The three headlines are then upgraded from ``some representative realises $L$'' to genuine
equalities of the descended function $\operatorname{flowExponentAt}$, the base generator $M$ having
determinant $1 \ne 0$ (so the descent is total).

\begin{theorem}[Representative-free cat exponent]\label{thm:susp-catFlowExpEq}
  \lean{ErgodicTheory.CatMapToral.catSuspension_flowExponentAt_eq_log}\leanok
  \uses{thm:susp-flowExpEq, thm:susp-catHasFlow}
  For $\hat\mu$-a.e.\ $q$, the descended flow exponent is
  $\operatorname{flowExponentAt} q = \log((3+\sqrt5)/2)$.
\end{theorem}
\begin{proof}\leanok
  \uses{thm:susp-flowExpEq, thm:susp-catHasFlow}
  Apply \Cref{thm:susp-flowExpEq} to the existential exponent \Cref{thm:susp-catHasFlow}; global
  invertibility of $M$ makes the lift total.
\end{proof}

\begin{theorem}[Representative-free Abramov reading]\label{thm:susp-catFlowExpBaseDiv}
  \lean{ErgodicTheory.CatMapToral.catSuspension_flowExponentAt_eq_base_div_roof}\leanok
  \uses{thm:susp-catFlowExpEq}
  The descended flow exponent equals the base top exponent of $M$ divided by $\int\tau$, an honest
  equality of the $\Sigma \to \R$ function rather than an existential over representatives.
\end{theorem}
\begin{proof}\leanok
  \uses{thm:susp-catFlowExpEq}
  Rewrite the base exponent as $\log((3+\sqrt5)/2)$ and $\int\tau = 1$ in
  \Cref{thm:susp-catFlowExpEq}.
\end{proof}

\begin{theorem}[Representative-free positivity]\label{thm:susp-catFlowExpPos}
  \lean{ErgodicTheory.CatMapToral.catSuspension_flowExponentAt_pos}\leanok
  \uses{thm:susp-catFlowExpEq}
  For $\hat\mu$-a.e.\ $q$, the descended flow exponent is strictly positive:
  $0 < \operatorname{flowExponentAt} q$.
\end{theorem}
\begin{proof}\leanok
  \uses{thm:susp-catFlowExpEq}
  Substitute the value $\log((3+\sqrt5)/2)$ from \Cref{thm:susp-catFlowExpEq}; it is positive
  because $(3+\sqrt5)/2 > 1$.
\end{proof}

\section{Ergodicity of the time-one map}

Fix a \emph{constant} roof $\tau \equiv r$. The time-$1$ map of the suspension flow acts on the
lifted plane by $(x, s) \mapsto (x, s + 1)$ modulo the deck relation, so its ergodicity is a
fibre-Fourier question with transform window the \emph{invariance period} $1$ --- not the roof $r$.
This is the constant-roof special-flow dichotomy of Cornfeld--Fomin--Sinai (Ch.~11). The spectral
input is that the base has no nontrivial unimodular eigenfunctions, which for a strongly mixing base
is automatic.

\begin{theorem}[Strong mixing of the two-sided Bernoulli shift]\label{thm:susp-bernMixing}
  \lean{ErgodicTheory.Multifractal.tendsto_measureReal_inter_preimage_iterate}\leanok
  For the invertible two-sided Bernoulli shift with i.i.d.\ measure $\operatorname{bernZ}\nu$ and
  measurable sets $A, B$, the diagonal correlations converge to the product of measures:
  \[
    \hat\nu(A \cap \sigma^{-k} B) \to \hat\nu(A)\,\hat\nu(B) \qquad (k \to \infty).
  \]
\end{theorem}
\begin{proof}\leanok
  Approximate $A, B$ by finite-block cylinder sets in symmetric difference; on cylinders the
  correlation is eventually exactly the product (independence past the block width), and the
  measure-preserving iterate keeps the symmetric-difference mass, so the $\varepsilon/5$
  bookkeeping pushes the exact identity to the limit. Classical (Cornfeld--Fomin--Sinai \S10;
  Walters, \emph{An Introduction to Ergodic Theory}, Thm.~1.30): Bernoulli shifts are strong mixing.
\end{proof}

\begin{theorem}[Mixing kills eigenvalues]\label{thm:susp-eigZero}
  \lean{ErgodicTheory.eigenfunction_ae_zero_of_mixing}\leanok
  A measurable eigenfunction $g$ (with $g \circ f = l\,g$) of a strongly-mixing transformation $f$,
  whose eigenvalue $l$ is unimodular ($\norm{l} = 1$) but $l \ne 1$, vanishes almost everywhere.
\end{theorem}
\begin{proof}\leanok
  Purely set-theoretic. The powers $l^n$ stay a fixed distance $\delta = \norm{l - 1}/2$ from $1$
  infinitely often (the two-consecutive-powers trick,
  \lean{ErgodicTheory.frequently_pow_far_from_one}). If $g$ were not a.e.\ zero, pick an admissible
  ball $B = \operatorname{ball} q\,\rho$ with $q \ne 0$ and $2\rho \le \norm{q}\delta$ whose
  preimage $A = g^{-1}B$ has positive measure. A point of $A \cap (f^n)^{-1}A$ would place both
  $g x$ and $l^n g x$ inside $B$, forcing $\norm{l^n q - q} < 2\rho \le \norm{q}\norm{l^n - 1}
  = \norm{l^n q - q}$ whenever $\norm{l^n - 1} \ge \delta$; this happens frequently, so the diagonal
  correlation vanishes along a subsequence --- contradicting mixing, which drives it to
  $\hat\mu(A)^2 > 0$.
\end{proof}

The harmonic-analysis engine works on the lifted plane. The transform window is the unit invariance
period, so no lap decomposition of the roof is needed.

\begin{definition}[The fibre Fourier coefficient]\label{thm:susp-coeffFn}
  \lean{ErgodicTheory.coeffFn}\leanok
  For $F : X \times \R \to \C$ the $n$-th \emph{fibre Fourier coefficient} over the unit invariance
  window is $\operatorname{coeffFn} F\,n\,x = \int_0^1 \overline{e_n(s)}\,F(x, s)\,ds$, where $e_n$
  is the $n$-th character of the length-$1$ circle. Its window is the period of the time-one map,
  not the roof.
\end{definition}

\begin{theorem}[The twisted eigenfunction relation]\label{thm:susp-twist}
  \lean{ErgodicTheory.coeffFn_twist}\leanok
  \uses{thm:susp-coeffFn}
  If $F$ is $1$-periodic in the fibre and satisfies the deck identity $F(Tx, s) = F(x, s + r)$, then
  \[
    \operatorname{coeffFn} F\,n\,(Tx) = e^{2\pi i n r}\,\operatorname{coeffFn} F\,n\,x.
  \]
\end{theorem}
\begin{proof}\leanok
  \uses{thm:susp-coeffFn}
  A three-step change of variables over the unit window: character algebra factoring out
  $e^{2\pi i n r}$, the translation $\int_0^1 G(s + r)\,ds = \int_r^{1+r} G$, and
  $1$-periodicity of $G$ to return the window to $[0, 1]$.
\end{proof}

For a $\zeta_1$-invariant set $A$ in a constant-roof suspension, its lifted indicator
$\operatorname{liftedIndicator} A = \mathbf 1_{\pi^{-1}A}$ (\lean{ErgodicTheory.liftedIndicator}) is
$1$-periodic in the fibre and satisfies the deck identity, so its fibre coefficients are twisted
eigenfunctions (\lean{ErgodicTheory.coeffFn_liftedIndicator_twist}).

\begin{theorem}[Per-fibre Parseval bridge]\label{thm:susp-parseval}
  \lean{ErgodicTheory.parseval_bridge}\leanok
  For a bounded measurable $f : \R \to \C$, the squared fibre-coefficient sum equals the fibre
  $L^2$ mass: $\sum_{n \in \Z}\norm{\operatorname{coeffFn}(f)\,n}^2 = \int_0^1\norm{f}^2$.
\end{theorem}
\begin{proof}\leanok
  Lift $f$ to the length-$1$ circle $\operatorname{AddCircle}(1)$; the interval coefficients are the
  circle's Fourier coefficients and the interval $L^2$ mass is the circle $L^2$ mass, so Mathlib's
  circle Parseval identity $\sum\norm{\hat f(n)}^2 = \norm{f}_2^2$ gives the claim.
\end{proof}

\begin{theorem}[Indicator dichotomy]\label{thm:susp-dichotomy}
  \lean{ErgodicTheory.fibre_indicator_dichotomy}\leanok
  \uses{thm:susp-parseval}
  Let $g : \R \to \R$ be measurable, $1$-periodic, $\{0, 1\}$-valued, with all nonzero-mode Fourier
  coefficients (over period $1$) vanishing. Then $g = 0$ a.e.\ on $(0, 1]$ or $g = 1$ a.e.\ on
  $(0, 1]$.
\end{theorem}
\begin{proof}\leanok
  \uses{thm:susp-parseval}
  Parseval collapses to the zero mode: $\bigl(\int_0^1 g\bigr)^2 = \int_0^1\norm{g}^2 = \int_0^1 g$
  (the last since $g^2 = g$ on $\{0, 1\}$), so $\int_0^1 g \in \{0, 1\}$. If the integral is $0$ then
  $g = 0$ a.e.; if it is $1$ then $1 - g = 0$ a.e. A periodic null-set spread
  (\lean{ErgodicTheory.measure_periodic_ae_zero_spread}) then carries the null fibre across every
  unit cell.
\end{proof}

\begin{theorem}[Time-one ergodicity, abstract base-generic form (issue \#35)]\label{thm:susp-timeOneErgodic}
  \lean{ErgodicTheory.ergodic_suspensionFlowMap_one_const_roof}\leanok
  \uses{thm:susp-twist, thm:susp-dichotomy, thm:susp-prob}
  Let $T$ be ergodic and measure-preserving on a probability space whose only measurable
  eigenfunction with a unimodular eigenvalue $\ne 1$ is $0$. For a positive \emph{irrational} roof
  $r$, the time-$1$ map of the constant-roof suspension flow is ergodic for $\hat\mu$.
\end{theorem}
\begin{proof}\leanok
  \uses{thm:susp-twist, thm:susp-dichotomy}
  Fix a $\zeta_1$-invariant $A$ with lifted indicator $F$. For $n \ne 0$ the twist eigenvalue
  $e^{2\pi i n r}$ is unimodular and $\ne 1$ (irrationality of $r$), so the spectral hypothesis
  forces $\operatorname{coeffFn} F\,n = 0$ a.e.\ (\Cref{thm:susp-twist}); the zero mode is
  $T$-invariant, hence a.e.\ constant by base ergodicity. Per fibre, the dichotomy
  \Cref{thm:susp-dichotomy} makes $F(x, \cdot)$ a.e.\ $0$ or a.e.\ $1$, and the periodic spread
  fills the roof window $[0, r)$; Fubini over the fundamental box then gives
  $\hat\mu(A) \in \{0, 1\}$.
\end{proof}

\begin{theorem}[Time-one ergodicity of the irrational-roof Bernoulli suspension]\label{thm:susp-bernTimeOne}
  \lean{ErgodicTheory.Multifractal.ergodic_suspensionFlow_timeOne_const_irrational}\leanok
  \uses{thm:susp-timeOneErgodic, thm:susp-bernMixing, thm:susp-eigZero}
  For the two-sided Bernoulli shift with measure $\operatorname{bernZ}\nu$ and any positive
  irrational roof $r$, the time-$1$ map of the constant-roof suspension flow is ergodic for the
  invariant suspension probability measure.
\end{theorem}
\begin{proof}\leanok
  \uses{thm:susp-timeOneErgodic, thm:susp-bernMixing, thm:susp-eigZero}
  Discharge the abstract \Cref{thm:susp-timeOneErgodic} with base ergodicity of the shift and the
  no-nontrivial-eigenfunctions input: strong mixing (\Cref{thm:susp-bernMixing}) feeds
  \Cref{thm:susp-eigZero}, forcing every measurable unimodular eigenfunction with eigenvalue $\ne 1$
  to vanish a.e.
\end{proof}

\begin{theorem}[Concrete non-vacuity witness $r = \sqrt2$]\label{thm:susp-sqrt2}
  \lean{ErgodicTheory.Multifractal.ergodic_suspensionFlow_timeOne_sqrtTwo}\leanok
  \uses{thm:susp-bernTimeOne}
  The irrational roof $r := \sqrt2$ yields an ergodic time-$1$ map, confirming the hypothesis
  $\operatorname{Irrational} r$ is satisfiable.
\end{theorem}
\begin{proof}\leanok
  \uses{thm:susp-bernTimeOne}
  Instantiate \Cref{thm:susp-bernTimeOne} at $r = \sqrt2$, positive and irrational.
\end{proof}

The irrationality is essential. For the \emph{unit} roof $r = 1$ the time-$1$ map is genuinely
non-ergodic (\Cref{ergodic_bernSuspensionFlow} records that only the full flow is ergodic).

\begin{theorem}[Non-ergodicity at the unit roof]\label{thm:susp-notErgodic}
  \lean{ErgodicTheory.Multifractal.not_ergodic_bernSuspensionFlow_one}\leanok
  \uses{bernSuspensionFlow}
  The time-$1$ map of the \emph{unit-roof} Bernoulli suspension flow is \emph{not} ergodic.
\end{theorem}
\begin{proof}\leanok
  \uses{bernSuspensionFlow}
  With $r = 1$ the deck translation advances the fibre by exactly the invariance period, so the
  saturated half-section $\{[x, s] : \operatorname{frac} s < 1/2\}$ is a nontrivial time-$1$-invariant
  set of mass $1/2$. Equivalently, the flow-generator eigenfunction $e^{2\pi i s}$ has eigenvalue
  $e^{2\pi i} = 1$ at time $1$ and descends to a non-constant invariant function; the zero-one law
  fails.
\end{proof}

\section{Entropy descent}

The Kolmogorov--Sinai entropy of the flow's time-$t$ maps is governed by the discrete power rule
and the Abramov time-change. The base algebra is the entropy analogue of \Cref{thm:susp-limitTransfer}.

\begin{theorem}[Discrete entropy power rule]\label{thm:susp-ksPow}
  \lean{ErgodicTheory.Entropy.ksEntropy_iterate}\leanok
  \uses{ksEntropy}
  For a measure-preserving $T$ on a probability space and $n \in \N$, $h(T^n) = n \cdot h(T)$.
\end{theorem}
\begin{proof}\leanok
  \uses{ksEntropy}
  Walters, \emph{An Introduction to Ergodic Theory}, Theorem 4.13. The $n$-fold refinement of a
  partition under $T^n$ is a subsequence of the refinement under $T$, and the Fekete limit along
  the arithmetic progression $kn$ equals $n$ times the full limit.
\end{proof}

\begin{theorem}[Flow iterate identity]\label{thm:susp-flowIter}
  \lean{ErgodicTheory.flow_iterate}\leanok
  For a measure-preserving flow $\varphi$, the $n$-th iterate of the time-$t$ map is the
  time-$(nt)$ map: $(\varphi_t)^{[n]} = \varphi_{nt}$.
\end{theorem}
\begin{proof}\leanok
  Induction on $n$ using $\varphi_{s+t} = \varphi_s \circ \varphi_t$ and the iterate recursion.
\end{proof}

\begin{theorem}[Flow entropy homogeneity along $\N$-multiples]\label{thm:susp-nsmulKs}
  \lean{ErgodicTheory.nsmul_ksEntropy_flow}\leanok
  \uses{thm:susp-ksPow, thm:susp-flowIter}
  For a measure-preserving flow $\varphi$ on a probability space,
  $n \cdot h(\varphi_t) = h(\varphi_{nt})$.
\end{theorem}
\begin{proof}\leanok
  \uses{thm:susp-ksPow, thm:susp-flowIter}
  Apply the power rule \Cref{thm:susp-ksPow} to $T = \varphi_t$, whose $n$-th iterate is
  $\varphi_{nt}$ by \Cref{thm:susp-flowIter}; entropy depends only on the underlying map.
\end{proof}

The Abramov time-change conjugates a constant-$r$ roof to the unit roof, rescaling the fibre.

\begin{definition}[The fibre-rescaling equivalence]\label{thm:susp-rescale}
  \lean{ErgodicTheory.suspensionRescale}\leanok
  \uses{thm:susp-space}
  The fibre map $(x, s) \mapsto (x, s/r)$ descends to a measurable equivalence
  $\operatorname{suspensionRescale} : \operatorname{SuspensionSpace} T\,(\tau \equiv r) \simeq
  \operatorname{SuspensionSpace} T\,(\tau \equiv 1)$, conjugating the constant-$r$ generator
  $(x, s) \mapsto (Tx, s - r)$ to the unit generator $(x, s) \mapsto (Tx, s - 1)$ (since
  $(s - r)/r = s/r - 1$). It is the constant-roof case of the Ambrose--Kakutani/Abramov time-change.
\end{definition}

\begin{theorem}[The time-$r$ entropy conjugacy]\label{thm:susp-rescaleKs}
  \lean{ErgodicTheory.ksEntropy_suspensionFlowMap_const_eq_unit}\leanok
  \uses{thm:susp-rescale}
  The time-$t$ map of the constant-$r$ suspension flow has the same Kolmogorov--Sinai entropy as the
  time-$(t/r)$ map of the unit-roof suspension flow.
\end{theorem}
\begin{proof}\leanok
  \uses{thm:susp-rescale}
  Measurable-conjugacy invariance of entropy applied to $\operatorname{suspensionRescale}$: it
  intertwines $\zeta^{(r)}_t$ with $\zeta^{(1)}_{t/r}$ and transports the invariant measure (the box
  $X \times [0, r)$ maps to $X \times [0, 1)$, scaling the fibre length by $r$, exactly cancelling
  the $r^{-1}$ vs $1^{-1}$ normalisations).
\end{proof}

\begin{theorem}[Time-$r$ entropy of the Bernoulli suspension]\label{thm:susp-timeRBern}
  \lean{ErgodicTheory.ksEntropy_bernConstSuspension_time_r}\leanok
  \uses{thm:susp-rescaleKs, ksEntropy_bernSuspensionFlow_one_eq_Hnu}
  For the constant-roof ($\tau \equiv r$) suspension of the two-sided Bernoulli shift, the time-$r$
  map has Kolmogorov--Sinai entropy exactly the per-symbol Shannon entropy $H_\nu$.
\end{theorem}
\begin{proof}\leanok
  \uses{thm:susp-rescaleKs}
  \Cref{thm:susp-rescaleKs} at $t = r$ (so $t/r = 1$) reduces to the unit-roof time-$1$ value
  $h(\zeta^{(1)}_1) = H_\nu$ (\Cref{ksEntropy_bernSuspensionFlow_one_eq_Hnu}).
\end{proof}

\begin{theorem}[Constant-roof time-one entropy, rational roof (issue \#38)]\label{thm:susp-timeOneEntropy}
  \lean{ErgodicTheory.ksEntropy_bernConstSuspension_time_one}\leanok
  \uses{thm:susp-nsmulKs, thm:susp-rescaleKs, ksEntropy_bernSuspensionFlow_one_eq_Hnu}
  For a \emph{rational} roof $r = a/b$ ($a, b \in \N$ positive), the time-$1$ map of the
  constant-roof Bernoulli suspension flow has Kolmogorov--Sinai entropy $H_\nu/r$.
\end{theorem}
\begin{proof}\leanok
  \uses{thm:susp-nsmulKs, thm:susp-rescaleKs}
  Fibre time-rescaling reduces $h(\zeta^{(r)}_1)$ to the unit-roof $h(\zeta^{(1)}_{1/r})$; writing
  $1/r = b/a$, homogeneity along $\N$-multiples (\Cref{thm:susp-nsmulKs}) gives
  $a \cdot h(\zeta^{(1)}_{b/a}) = h(\zeta^{(1)}_b) = b \cdot h(\zeta^{(1)}_1) = b\,H_\nu$. Since
  $H_\nu$ is finite and $a > 0$, the $\N$-scalar $\overline\R$-equation $a \cdot h(\zeta^{(r)}_1)
  = b\,H_\nu$ pins $h(\zeta^{(r)}_1) = b\,H_\nu/a = H_\nu/r$, the arithmetic descending to $\R$.
\end{proof}

\begin{theorem}[Multiplicative form]\label{thm:susp-timeOneEntropyMul}
  \lean{ErgodicTheory.ksEntropy_bernConstSuspension_time_one_mul}\leanok
  \uses{thm:susp-timeOneEntropy}
  For a rational roof $r = a/b$, $\;h(\zeta^{(r)}_1) \cdot r = H_\nu$.
\end{theorem}
\begin{proof}\leanok
  \uses{thm:susp-timeOneEntropy}
  Multiply the value form \Cref{thm:susp-timeOneEntropy} by $r$; the product
  $(H_\nu/r) \cdot r = H_\nu$ is real arithmetic.
\end{proof}

\section{Abstract Abramov flow-entropy homogeneity}

The rational restriction of \Cref{thm:susp-timeOneEntropy} is an artefact of the discrete power
rule, not of the mathematics: the full Abramov homogeneity $h(\varphi_t) = t\,h(\varphi_1)$ holds
for \emph{every} real time $t > 0$. We prove it abstractly for any measure-continuous
measure-preserving flow on a standard Borel probability space, following Ito's elementary
generator-free argument (Y.~Ito, \emph{An elementary proof of Abramov's result on the entropy of a
flow}, Nagoya Math.\ J.\ \textbf{41} (1971), 1--5), which avoids the flow generators / Rokhlin
towers (Ambrose--Kakutani structure) of Abramov's original approach (L.~M.~Abramov, \emph{On the
entropy of a flow}, Dokl.\ Akad.\ Nauk SSSR \textbf{128} (1959) 873--875).

\begin{definition}[Measure-continuity of a flow]\label{def:flow-measCont}
  \lean{ErgodicTheory.MeasurePreservingFlow.MeasureContinuous}\leanok
  A measure-preserving flow $\varphi$ is \emph{measure-continuous} if for every measurable set $A$
  the real measure of the symmetric difference $\varphi_t^{-1} A \,\triangle\, A$ tends to $0$ as
  $t \to 0$.
\end{definition}

\begin{theorem}[Measure-continuity of the suspension flow]\label{thm:flow-suspMeasCont}
  \lean{ErgodicTheory.tendsto_measureReal_symmDiff_suspensionFlowMap}\leanok
  \uses{bernSuspensionFlow}
  The unit-roof suspension flow of an \emph{arbitrary} measure-preserving base is measure-continuous:
  for every measurable $A$, $\mu\!\left(\varphi_t^{-1} A \,\triangle\, A\right) \to 0$ as $t \to 0$.
  This is the keystone hypothesis; the Bernoulli suspension flow instantiates it
  (\lean{ErgodicTheory.measureContinuous_bernSuspensionFlow}).
\end{theorem}
\begin{proof}\leanok
  Fibre-translation continuity of the roof-1 fundamental box, combined with dominated convergence
  over the box.
\end{proof}

\begin{theorem}[Ito's L1: the partition moves little under a small shift]\label{thm:flow-L1}
  \lean{ErgodicTheory.condEntropyGivenPartition_flow_tendsto_zero}\leanok
  \uses{def:flow-measCont}
  For a measure-continuous flow and a finite partition $P$, the conditional Shannon entropy
  $H(\varphi_t P \mid P) \to 0$ as $t \to 0$ (Ito's (2.2)).
\end{theorem}
\begin{proof}\leanok
  \uses{def:flow-measCont}
  The self-conditioning $H(P \mid P) = 0$, and the cell measures $\mu(\varphi_t^{-1}P_i \cap P_j)$
  are continuous at $t = 0$ by measure-continuity; continuity of the Shannon cell functional
  transports the limit.
\end{proof}

\begin{theorem}[Two-family Shannon comparison]\label{thm:flow-finJoin}
  \lean{ErgodicTheory.Entropy.entropy_finJoin_le_add_sum_condEntropy}\leanok
  For two finite families $\beta = (\beta_k)$, $\gamma = (\gamma_k)$ of finite partitions,
  $H\!\left(\bigvee_k \gamma_k\right) \le H\!\left(\bigvee_k \beta_k\right)
  + \sum_k H(\gamma_k \mid \beta_k)$.
\end{theorem}
\begin{proof}\leanok
  Pure Shannon entropy: refinement monotonicity, the chain rule, conditional subadditivity over the
  family, and anti-monotonicity of conditioning.
\end{proof}

\begin{theorem}[$\varepsilon$--$\delta$ alignment proposition]\label{thm:flow-ratioLUB}
  \lean{ErgodicTheory.exists_isLUB_ksEntropyPartition_flow_ratio}\leanok
  \uses{def:flow-measCont, thm:flow-L1, thm:flow-finJoin}
  For a measure-continuous flow on a standard Borel space and a finite partition $P$, the ratio
  $t \mapsto h(\varphi_t, P)/t$ converges as $t \downarrow 0$ to its least upper bound
  $L \in \overline{\R}$ over all positive times. The statement is genuinely in $\overline{\R}$: the
  $\R$-valued form is false for infinite-entropy flows, and this repair is disclosed in-module.
\end{theorem}
\begin{proof}\leanok
  \uses{thm:flow-L1, thm:flow-finJoin}
  The alignment inequality (from \Cref{thm:flow-L1} and \Cref{thm:flow-finJoin}) makes each slope
  eventually dominate every value below the supremum; the supremum is thus the limit.
\end{proof}

\begin{theorem}[Abstract Abramov homogeneity]\label{thm:flow-abramov-hom}
  \lean{ErgodicTheory.ksEntropy_flow_eq_mul}\leanok
  \uses{thm:flow-ratioLUB, thm:susp-nsmulKs}
  For a measure-continuous measure-preserving flow $\varphi$ on a standard Borel probability space
  and $t > 0$, $\;h(\varphi_t) = t \cdot h(\varphi_1)$ in $\overline{\R}$.
\end{theorem}
\begin{proof}\leanok
  \uses{thm:flow-ratioLUB, thm:susp-nsmulKs}
  Interchange the supremum over partitions with the alignment proposition
  (\Cref{thm:flow-ratioLUB}), using the discrete $\N$-homogeneity (\Cref{thm:susp-nsmulKs}); at
  $t = 1$ the slope is $h(\varphi_1)$, giving the stated multiple.
\end{proof}

\begin{theorem}[Unit-roof time-$s$ entropy, all $s > 0$]\label{thm:flow-timeS}
  \lean{ErgodicTheory.ksEntropy_bernSuspensionFlow_time_s_eq}\leanok
  \uses{thm:flow-abramov-hom, ksEntropy_bernSuspensionFlow_one_eq_Hnu}
  For the unit-roof Bernoulli suspension flow and every $s > 0$,
  $h(\zeta^{(1)}_s) = s \cdot H_\nu$.
\end{theorem}
\begin{proof}\leanok
  \uses{thm:flow-abramov-hom}
  \Cref{thm:flow-abramov-hom} for the measure-continuous Bernoulli suspension flow, at the finite
  value $h(\zeta^{(1)}_1) = H_\nu$.
\end{proof}

\begin{theorem}[Constant-roof time-one entropy, all roofs (issue \#48)]\label{thm:flow-timeOneAll}
  \lean{ErgodicTheory.ksEntropy_bernConstSuspension_time_one_irrational}\leanok
  \uses{thm:flow-timeS}
  For \emph{every} roof $r > 0$ (irrational included), the time-$1$ map of the constant-roof
  Bernoulli suspension flow has Kolmogorov--Sinai entropy $H_\nu/r$.
\end{theorem}
\begin{proof}\leanok
  \uses{thm:flow-timeS}
  Fibre time-rescaling reduces $h(\zeta^{(r)}_1)$ to the unit-roof $h(\zeta^{(1)}_{1/r})$;
  \Cref{thm:flow-timeS} at $s = 1/r$ gives $(1/r)\,H_\nu = H_\nu/r$.
\end{proof}

This retires the entropy-side wall of \Cref{thm:susp-timeOneEntropy}: the value $H_\nu/r$ now holds,
with a formalized proof, for \emph{every} roof $r > 0$. The two headline features of the constant-roof
time-$1$ map are then \emph{complementary}. Ergodicity holds exactly for \emph{irrational} roofs
(\Cref{thm:susp-bernTimeOne}) and fails for rational ones (\Cref{thm:susp-notErgodic}); the entropy
is $H_\nu/r$ for all $r > 0$ (\Cref{thm:flow-timeOneAll}), so on the irrational roofs a single object
carries both an ergodic time-$1$ map and pinned positive entropy $H_\nu/r$. The ergodicity turns on
the deck twist $e^{2\pi i n r}$ being nontrivial (irrational $r$); the entropy, on the flow
homogeneity, which no longer sees the arithmetic of $r$.

\section{The quotient-level suspension flow cocycle}\label{sec:susp-quotientCocycle}

The representative-free exponent above lives on orbit classes but is built from a chosen
representative. Issue \#53 asks for the strictly finer object: the flow's \emph{own} matrix
derivative data, packaged as a genuine \Cref{FlowCocycle} on the quotient (mapping-torus) space, the
very interface consumed by the continuous-flow MET \Cref{oseledets_flow}. This is achievable for the
\textbf{constant unit roof} $\tau \equiv 1$, where the fundamental domain is the box $X \times [0,1)$
and the quotient projection restricted to it is a \emph{measurable bijection} --- the measurable
trivialization of the suspension bundle over a standard Borel base.

\begin{definition}[The two-sided $\Z$-indexed matrix cocycle]\label{cocycleZ}
  \lean{ErgodicTheory.cocycleZ}\leanok
  \uses{FlowCocycle}
  Over an invertible measure-preserving base $T : X \simeq X$, the one-sided iterated cocycle
  (indexed by $\N$) extends to a two-sided $\Z$-indexed cocycle $\operatorname{cocycleZ} A\,T\,n$: for
  $n \ge 0$ the forward cocycle, and for $n < 0$ the backward cocycle of the inverse dynamics with
  generator $x \mapsto (A(T^{-1}x))^{-1}$.
\end{definition}

\begin{theorem}[The two-sided cocycle identity]\label{cocycleZ_add}
  \lean{ErgodicTheory.cocycleZ_add}\leanok
  \uses{cocycleZ}
  $\operatorname{cocycleZ}(m + n)\,x = \operatorname{cocycleZ} m\,(\operatorname{baseIter} n\,x)
  \cdot \operatorname{cocycleZ} n\,x$, where $\operatorname{baseIter}$ is the $\Z$-iterate of the
  base. This is the keystone that lets the discrete object assemble into a continuous-time cocycle.
\end{theorem}

\begin{definition}[The genuine quotient flow cocycle]\label{quotientFlowCocycle}
  \lean{ErgodicTheory.quotientFlowCocycle}\leanok
  \uses{cocycleZ, FlowCocycle}
  For the constant unit roof, reading off a measurable canonical representative
  $\operatorname{unitFwd} q = (\operatorname{baseIter}\lfloor s\rfloor x, \{s\})$ of each orbit class
  $q = [x,s]$, the assignment $t \mapsto \operatorname{cocycleZ} A\,T\,\lfloor \{s\} + t\rfloor\,
  (\operatorname{unitFwd} q)_1$ is a genuine \Cref{FlowCocycle} on the mapping torus. The four cocycle
  fields follow from \Cref{cocycleZ_add} transported through the floor split
  $\lfloor a + t\rfloor = \lfloor a\rfloor + \lfloor \{a\} + t\rfloor$.
\end{definition}

\begin{theorem}[Cohomology to the cover cocycle]\label{exists_flowCocycle_cohomologous_to_cover}
  \lean{ErgodicTheory.exists_flowCocycle_cohomologous_to_cover}\leanok
  \uses{quotientFlowCocycle, thm:susp-coverCocycle}
  Over the quotient, \Cref{quotientFlowCocycle} is cohomologous to the cover cocycle
  (\Cref{thm:susp-coverCocycle}): at a general representative $(x,s)$ with $0 \le s$ the two differ
  by the measurable \emph{rep-level frame} $C(x,s) = \operatorname{cocycleZ} A\,T\,\lfloor s\rfloor
  x$. Because different representatives of one class carry different $\lfloor s\rfloor$, the issue's
  schematic \emph{class}-level conjugacy is provably impossible; the honest general statement is this
  rep-level frame, collapsing to the conjugation-free canonical identity at $\operatorname{unitFwd}$.
\end{theorem}
\begin{proof}\leanok
  \uses{quotientFlowCocycle, cocycleZ_add}
  Take $B = \operatorname{quotientFlowCocycle}$ and $C(p) = \operatorname{cocycleZ} A\,T
  \lfloor p_2\rfloor p_1$; measurability of $C$ is the composition of $\lfloor \cdot\rfloor$ with the
  measurable-in-index cocycle, and the frame identity is the floor split of \Cref{cocycleZ_add}.
\end{proof}

\begin{theorem}[Exponent transport]\label{flowExponentAt_quotientFlowCocycle}
  \lean{ErgodicTheory.flowExponentAt_quotientFlowCocycle}\leanok
  \uses{quotientFlowCocycle, thm:susp-flowExpAt}
  Along the $\operatorname{atTop}$ half-line $0 \le t$, the growth rate of \Cref{quotientFlowCocycle}
  equals the descended flow exponent $\operatorname{flowExponentAt}$ (\Cref{thm:susp-flowExpAt}).
\end{theorem}

The general non-constant roof is \textbf{out of scope} and disclosed: the orbit quotient has no
measurable canonical representative there, so only the exponent, not the matrix cocycle, descends.

\subsection{The cat-map instance}

\begin{definition}[The cat map's quotient flow cocycle]\label{catQuotientFlowCocycle}
  \lean{ErgodicTheory.CatMapToral.catQuotientFlowCocycle}\leanok
  \uses{quotientFlowCocycle, catTorus}
  Instantiating \Cref{quotientFlowCocycle} at the Arnold cat map \Cref{catTorus} under the constant
  unit roof, with base generator the cat map's own derivative cocycle (the constant hyperbolic
  matrix $\operatorname{cat}_\R$), yields a genuine \Cref{FlowCocycle} on the cat suspension.
\end{definition}

\begin{theorem}[The cat quotient flow-cocycle Lyapunov exponent]\label{catQuotientFlowCocycle_exponent}
  \lean{ErgodicTheory.CatMapToral.catQuotientFlowCocycle_exponent}\leanok
  \uses{catQuotientFlowCocycle, flowExponentAt_quotientFlowCocycle, thm:susp-catFlowExpEq}
  For $\hat\mu$-a.e.\ orbit class $q$, the growth rate of \Cref{catQuotientFlowCocycle} converges to
  the flow Lyapunov exponent $\log\!\big((3+\sqrt5)/2\big)$ --- the log of the cat matrix's top
  eigenvalue.
\end{theorem}
\begin{proof}\leanok
  \uses{flowExponentAt_quotientFlowCocycle, thm:susp-catFlowExpEq}
  Chain the descended cat-suspension exponent (\Cref{thm:susp-catFlowExpEq}) through the
  exponent-transport \Cref{flowExponentAt_quotientFlowCocycle}.
\end{proof}
